\section{Quantum Logic Gates and Circuits}
Just as classical algorithms rely on networks of fundamental Boolean logic gates like AND, OR, and NOT, quantum algorithms are built using the \textit{quantum circuit model}.
In this model, computations are performed by applying a carefully timed sequence of operations—known as quantum logic gates—to a register of initialized qubits.
The computation then concludes with a classical measurement to extract the final deterministic result \cite{deutsch1989quantum}.

A core principle of closed quantum systems is that their evolution over time is both deterministic and entirely reversible.
Because of this, every single quantum logic gate must be a mathematically invertible operation.
The only exception to this rule is the final measurement, which irreversibly collapses the delicate quantum state.
Additionally, because these gates manipulate probability amplitudes, they must preserve the fundamental rule that all probabilities sum exactly to 100\%.
To satisfy these twin requirements of reversibility and probability conservation, all quantum gates are mathematically represented as \textit{unitary matrices} \cite{nielsen2010quantum}.

\subsection{Single-Qubit Gates}
Single-qubit gates operate on the state space of an individual qubit, $\mathbb{C}^2$, and are represented by $2 \times 2$ unitary matrices.
The most foundational of these are the three Pauli matrices.
Geometrically, these correspond to 180-degree ($\pi$ radian) rotations around the $x$, $y$, and $z$ axes of the Bloch sphere:
\begin{align}
    X &= \begin{pmatrix} 0 & 1 \\ 1 & 0 \end{pmatrix} \\
    Y &= \begin{pmatrix} 0 & -i \\ i & 0 \end{pmatrix} \\
    Z &= \begin{pmatrix} 1 & 0 \\ 0 & -1 \end{pmatrix}
\end{align}
The $X$ gate is the quantum equivalent of the classical NOT gate; it acts as a bit-flip, mapping $\ket{0} \to \ket{1}$ and $\ket{1} \to \ket{0}$.
The $Z$ gate, on the other hand, acts as a phase-flip.
It leaves the $\ket{0}$ state completely unchanged but flips the sign of the $\ket{1}$ state, mapping $\ket{1} \to -\ket{1}$.

Perhaps the most ubiquitous single-qubit operation in all of quantum computing is the Hadamard gate, denoted $H$.
This gate is the primary tool for generating quantum superposition.
When applied to a deterministic computational basis state, it creates a perfectly equal probability distribution between $\ket{0}$ and $\ket{1}$:
\begin{equation}
    H = \frac{1}{\sqrt{2}}\begin{pmatrix} 1 & 1 \\ 1 & -1 \end{pmatrix}
\end{equation}
Specifically, applying $H$ to $\ket{0}$ yields the positive superposition state $\ket{+} = \frac{1}{\sqrt{2}}(\ket{0} + \ket{1})$.
Applying it to $\ket{1}$ yields the negative superposition state $\ket{-} = \frac{1}{\sqrt{2}}(\ket{0} - \ket{1})$ \cite{mermin2007quantum}.

Other essential single-qubit operations include the Phase gate ($S$) and the $\pi/8$ gate ($T$).
Instead of flipping the state, these gates gently rotate it around the Z-axis, introducing important relative phases without altering the measurement probabilities:
\begin{equation}
    S = \begin{pmatrix} 1 & 0 \\ 0 & i \end{pmatrix}, \quad T = \begin{pmatrix} 1 & 0 \\ 0 & e^{i\pi/4} \end{pmatrix}
\end{equation}

\subsection{Two-Qubit Gates: CNOT, SWAP, and Controlled-Phase}
While single-qubit gates can explore individual state spaces, we must apply gates across multiple qubits to perform conditional logic and generate entanglement.
The quintessential two-qubit operation is the Controlled-NOT (CNOT) gate.

The CNOT gate involves two distinct roles: a control qubit and a target qubit.
If the control qubit is in the $\ket{1}$ state, the gate applies an $X$ gate (a bit-flip) to the target qubit.
If the control qubit is in the $\ket{0}$ state, the target qubit is left completely untouched.
In the standard computational basis $\{\ket{00}, \ket{01}, \ket{10}, \ket{11}\}$, where the first qubit acts as the control, the matrix representation is:
\begin{equation}
    \text{CNOT} = \begin{pmatrix} 
    1 & 0 & 0 & 0 \\ 
    0 & 1 & 0 & 0 \\ 
    0 & 0 & 0 & 1 \\ 
    0 & 0 & 1 & 0 
    \end{pmatrix}
\end{equation}
The CNOT gate is absolutely critical because it serves as the primary mechanism for entangling previously independent qubits.
For example, if we start with the state $\ket{00}$, apply a Hadamard gate to the first qubit to create a superposition, and then apply a CNOT gate, we generate the maximally entangled Bell state $\ket{\Phi^+} = \frac{1}{\sqrt{2}}(\ket{00} + \ket{11})$ \cite{nielsen2010quantum}.

Another highly practical two-qubit operation is the SWAP gate, which simply exchanges the quantum states of two distinct qubits.
This gate is indispensable for routing information across physical hardware topologies where not all qubits are directly connected to each other.
It also forms the final bit-reversal step of the Quantum Fourier Transform (QFT) used in Shor's algorithm:
\begin{equation}
    \text{SWAP} = \begin{pmatrix} 
    1 & 0 & 0 & 0 \\ 
    0 & 0 & 1 & 0 \\ 
    0 & 1 & 0 & 0 \\ 
    0 & 0 & 0 & 1 
    \end{pmatrix}
\end{equation}

While the CNOT conditionally applies a bit-flip, the Controlled-Phase gate—denoted $CP(\theta)$ or $CR_z(\theta)$—applies a conditional phase shift.
This gate adds a complex phase of $e^{i\theta}$ to the system if, and only if, both the control and target qubits are simultaneously in the $\ket{1}$ state:
\begin{equation}
    CP(\theta) = \begin{pmatrix} 
    1 & 0 & 0 & 0 \\ 
    0 & 1 & 0 & 0 \\ 
    0 & 0 & 1 & 0 \\ 
    0 & 0 & 0 & e^{i\theta} 
    \end{pmatrix}
\end{equation}
The $CP$ gate acts as the mathematical cornerstone of the QFT.
It enables the precise extraction of periodic phases, a mechanism that is strictly necessary for tasks like integer factorization.

\subsection{Advanced and Multi-Controlled Gates}
To implement more complex logical constructs—like the Search Oracles in Grover's algorithm or the modular arithmetic in Shor's algorithm—we need gates that condition their actions on several qubits at once.

The Toffoli gate, also known as the Controlled-Controlled-NOT (CCNOT), is the most famous example of this.
It watches two control qubits and applies an $X$ gate to a single target qubit if, and only if, both controls are in the $\ket{1}$ state.
Because the Toffoli gate is proven to be universal for classical reversible computation, it is the key ingredient for building quantum oracles that evaluate classical boolean functions \cite{nielsen2010quantum}:
\begin{equation}
    \text{CCNOT} = \begin{pmatrix} 
    1 & 0 & 0 & 0 & 0 & 0 & 0 & 0 \\
    0 & 1 & 0 & 0 & 0 & 0 & 0 & 0 \\
    0 & 0 & 1 & 0 & 0 & 0 & 0 & 0 \\
    0 & 0 & 0 & 1 & 0 & 0 & 0 & 0 \\
    0 & 0 & 0 & 0 & 1 & 0 & 0 & 0 \\
    0 & 0 & 0 & 0 & 0 & 1 & 0 & 0 \\
    0 & 0 & 0 & 0 & 0 & 0 & 0 & 1 \\
    0 & 0 & 0 & 0 & 0 & 0 & 1 & 0 
    \end{pmatrix}
\end{equation}

For more advanced quantum algorithms, the Toffoli gate is often generalized into the Multi-Controlled X gate ($MCX$).
The $MCX$ gate flips a target qubit conditioned on the state of an arbitrary number of $n$ control qubits.

Similarly, the diffusion step in Grover's algorithm relies heavily on a Multi-Controlled Z gate ($MCZ$).
Instead of flipping the bit, the $MCZ$ gate flips the phase (the sign) of the entire system, but only when all $n$ qubits are simultaneously in the $\ket{1}$ state.

In practical hardware, we typically do not have native $MCZ$ gates available to us.
Instead, we can easily synthesize one by taking an $MCX$ gate and wrapping its target qubit in Hadamard gates.
This mathematical trick smoothly shifts the bit-flip operation into a phase-flip operation: 
$MCZ = (I \otimes \dots \otimes H) MCX (I \otimes \dots \otimes H)$ \cite{barenco1995elementary}.

Finally, many advanced algorithms rely on the generalized Controlled-$U$ operation ($c-U$).
Imagine an arbitrary quantum operation $U$ designed to act on a large target register of $n$ qubits.
The $c-U$ gate evaluates whether a single control qubit is $\ket{1}$, and if so, it triggers the entire $U$ operation on the target register.
This conditional triggering is the core engine behind the Quantum Phase Estimation subroutine.
It is exactly what allows algorithms to evaluate complex functions, such as modular exponentiation $a^x \pmod N$, in massive parallel across a quantum superposition.